\documentclass[aspectratio=169,11pt]{beamer}

% ============================================
% PACKAGES
% ============================================
\usepackage[utf8]{inputenc}
\usepackage[T1]{fontenc}
\usepackage[ngerman]{babel}
\usepackage{lmodern}
\usepackage{tcolorbox}
\tcbuselibrary{skins, hooks}
\usepackage{listings}
\usepackage{xcolor}
\usepackage{fontawesome5}

% Modern Font
\IfFileExists{FiraSans.sty}{\usepackage[sfdefault]{FiraSans}}{}
\renewcommand{\familydefault}{\sfdefault}

% ============================================
% COLORS & DESIGN SYSTEM
% ============================================
\definecolor{BgColor}{HTML}{FAFAFA}
\definecolor{Primary}{HTML}{2B3A42}
\definecolor{Accent}{HTML}{E84A5F}
\definecolor{Secondary}{HTML}{3F5765}
\definecolor{CodeBg}{HTML}{F0F0F0}
\definecolor{Success}{HTML}{28A745}

\setbeamercolor{background canvas}{bg=BgColor}
\setbeamercolor{title}{fg=Primary}
\setbeamercolor{frametitle}{bg=Primary, fg=white}
\setbeamercolor{structure}{fg=Accent}
\setbeamercolor{normal text}{fg=Primary}
\setbeamercolor{itemize item}{fg=Accent}

\setbeamertemplate{navigation symbols}{}
\setbeamertemplate{footline}{%
  \leavevmode%
  \hbox{%
  \begin{beamercolorbox}[wd=.333333\paperwidth,ht=2.25ex,dp=1ex,center]{author in head/foot}%
    \usebeamerfont{author in head/foot}\tiny Falko Himmel, Malte Baumann, Max Lehm
  \end{beamercolorbox}%
  \begin{beamercolorbox}[wd=.333333\paperwidth,ht=2.25ex,dp=1ex,center]{title in head/foot}%
    \usebeamerfont{title in head/foot}\tiny Grundlagen der Praktischen Informatik
  \end{beamercolorbox}%
  \begin{beamercolorbox}[wd=.333333\paperwidth,ht=2.25ex,dp=1ex,right]{date in head/foot}%
    \usebeamerfont{date in head/foot}\insertframenumber{} / \inserttotalframenumber\hspace*{2ex} 
  \end{beamercolorbox}}%
  \vskip0pt%
}

% ============================================
% CUSTOM BOXES
% ============================================
\tcbset{
    modernbox/.style={
        enhanced,
        boxrule=0pt,
        frame hidden,
        colback=white,
        coltitle=white,
        colbacktitle=Secondary,
        fonttitle=\bfseries,
        drop lifted shadow,
        arc=4pt,
        toptitle=2mm,
        bottomtitle=2mm,
        left=2mm,
        right=2mm
    },
    alertbox/.style={
        modernbox,
        colbacktitle=Accent
    }
}

% ============================================
% CODE LISTINGS
% ============================================
\lstdefinelanguage{Haskell}{
  keywords={class, data, deriving, do, else, if, in, import, let, module, of, then, type, where, case, error},
  keywordstyle=\color{Accent}\bfseries,
  identifierstyle=\color{Primary},
  sensitive=true,
  comment=[l]{--},
  commentstyle=\color{gray}\ttfamily\itshape,
  stringstyle=\color{teal}\ttfamily,
  morestring=[b]"
}

\lstset{
    language=Haskell,
    backgroundcolor=\color{CodeBg},
    basicstyle=\ttfamily\scriptsize,
    breaklines=true,
    frame=single,
    rulecolor=\color{CodeBg},
    frameround=tttt,
    numbers=left,
    numberstyle=\tiny\color{gray},
    xleftmargin=15pt,
    tabsize=4,
    showstringspaces=false
}

% ============================================
% TITLE PAGE
% ============================================
\title{\textbf{Tutorium 13}}
\subtitle{Tests, Beweise \& Altklausuraufgaben}
\institute{Grundlagen der Praktischen Informatik}
\author{}
\date{\today}

\begin{document}

\begin{frame}[plain]
    \titlepage
\end{frame}

\begin{frame}{Inhaltsverzeichnis}
    \tableofcontents
\end{frame}

% ============================================
% TESTS
% ============================================
\section{Tests}

\begin{frame}{Softwaretests}
    Um sicherzustellen, dass unser komplexer Algorithmus funktioniert, schreiben wir Tests.
    
    \begin{columns}[T]
        \begin{column}{0.48\textwidth}
            \begin{tcolorbox}[modernbox, title=\faBox\ Blackbox-Tests]
                \begin{itemize}
                    \item Prüfen nur Eingabe $\rightarrow$ Ausgabe
                    \item Überprüfen die Korrektheit gemäß Spezifikation
                    \item In Haskell nutzen wir oft \textbf{QuickCheck}
                \end{itemize}
            \end{tcolorbox}
        \end{column}
        \begin{column}{0.48\textwidth}
            \begin{tcolorbox}[modernbox, title=\faSearch\ Whitebox-Tests]
                \begin{itemize}
                    \item Prüfen die internen Abläufe (Verzweigungen)
                    \item Ziel: Hohe Testabdeckung (\textit{Code Coverage})
                    \item Eher seltener in der Klausur
                \end{itemize}
            \end{tcolorbox}
        \end{column}
    \end{columns}
\end{frame}

% ============================================
% BEWEISE
% ============================================
\section{Beweise (Induktion)}

\begin{frame}[fragile]{Korrektheit durch Induktion}
    In Haskell sind Algorithmen rekursiv definiert. Deshalb nutzen wir \textbf{vollständige Induktion}, um Eigenschaften mathematisch zu beweisen.
    
    \begin{tcolorbox}[modernbox, title=Beispiel: Beweis von \texttt{reverse (x:xs)}]
        \textbf{Basisfall:} \texttt{reverse [] = []} \quad (wahr laut Definition) \\
        \vspace{0.2cm}
        \textbf{Induktionsannahme (IA):} Angenommen, für eine beliebige Liste \texttt{xs} gilt: \\
        \texttt{reverse xs = rev xs} \\
        \vspace{0.2cm}
        \textbf{Induktionsschritt (IS):} Zeige, dass die Aussage auch für \texttt{(x:xs)} gilt. \\
        \begin{tabular}{ll}
            \texttt{reverse (x:xs)} & $= \texttt{reverse xs ++ [x]}$ \quad (Def. von reverse) \\
            & $= \texttt{rev xs ++ [x]}$ \quad (nach IA) \\
            & $= \texttt{umgedreht (x:xs)}$ \quad (Def. von rev) \\
        \end{tabular}
        \flushright \textit{q.e.d.}
    \end{tcolorbox}
\end{frame}

% ============================================
% LIVE DEMO
% ============================================
\section{Praxis}

\begin{frame}[plain]
    \begin{center}
        \Huge \textbf{Live Demo} \\
        \vspace{0.5cm}
        \large \faLaptopCode~ \texttt{Aufgabe3.hs} \\
        \vspace{0.3cm}
        \normalsize Wir programmieren weitere Altklausuraufgaben (\texttt{eitherFilter}, \texttt{list2multiset}, \texttt{selects}) live!
\vspace{0.5cm}\begin{center}\href{https://moodle.uni-ulm.de/pluginfile.php/1348606/mod_folder/content/0/2024_1_exam.pdf}{\beamerbutton{Altklausur 2024\_1}}\quad\href{https://moodle.uni-ulm.de/pluginfile.php/1348606/mod_folder/content/0/2024_2_exam.pdf}{\beamerbutton{Altklausur 2024\_2}}\end{center}
    \end{center}
\end{frame}

\end{document}
